\section{Eingriffsauflösung und Kopfkante}
\label{sec:kopfkante}

Die Stellungen des Eingriffs werden als getrennte statische Rechnungen über dem Eingriffsweg aufgereiht (vgl. \secref{sec:fe_modell}). Wie dicht sie liegen müssen, wurde mit 60 Stellungen je Teilung am Ringkörper unter W1 bestimmt. Von 133 Stellungen liefen 122 durch, elf brachen ab. Bei $16$\,Nm trägt der nachfolgende Zahn durch die Verformung bereits $1{,}65$\,mm vor A, und die Kopfkante des Ritzels bleibt am auslaufenden Zahn $1{,}5$\,mm hinter E im Kontakt. Die größte Fußspannung des Rades liegt deshalb im Bereich des Normpunktes C und nicht bei B. Bei $16$\,Nm halten 30 Stellungen je Teilung die Größtwerte innerhalb $0{,}15$\,\% und 20 innerhalb $1{,}6$\,\%, letzteres an einem Knick von $|\sigma_{\mathrm{3}}|$. Mit 20 je Teilung wandern die Maxima von Verformung und Fußspannung um weniger als $0{,}05$\,mm. Bei $8$\,Nm sind die Maxima Knicke an den Übergängen der Zahl tragender Paare. Dort verfehlt das Gitter mit 30 je Teilung einen Größtwert um $1{,}7$\,\%, das Gitter mit 20 trifft ihn.

Für die Serie aller vier Werkstoffstufen wurden 30 Stellungen je Teilung und ein Rand von einer Teilung vor A und hinter E gewählt. Ein Rand von einer halben Teilung war in einem Vorlauf mit W4 zu kurz, die äußeren Zähne trugen an den Gitterenden noch $10$\,\% und $30$\,\% der Zahnkraft. In der Serie trägt der mittlere Zahn bis zu $0{,}60$ Teilungen vor A und bis zu $0{,}48$ Teilungen hinter E, mit Moment und Plastizität wachsend. Innerhalb dieses Gitters verändern 15 Stellungen je Teilung die Größtwerte um höchstens $2{,}9$\,\% bei W3 und $16$\,Nm, 10 je Teilung um $5{,}5$\,\%.

Alle elf Abbrüche der Studie liegen an drei Wegstellen je Teilung. Dort sitzt die scharfe Kopfkante des starren Ritzels auf einer Knotenreihe der Flanke des auslaufenden Zahns hinter E. Die Knotenreihen liegen auf den Radien $25{,}698$, $25{,}800$ und $25{,}901$\,mm der Radflanke, also $0{,}1$\,mm auseinander. Der Knoten auf dem Knick meldet einen zu großen Durchdringungsfehler, das Inkrement wird bis auf $10^{-6}$ verkürzt, und die Rechnung endet. Weder eine Strafsteifigkeit statt der Lagrange-Multiplikatoren noch erhöhte Iterationsgrenzen helfen. Ursache ist die Mehrdeutigkeit der Flächennormalen an einer scharfen Ecke, der Löser wählt einen der beiden Zweige zufällig. Eine Rundung von $0{,}05$\,mm, kleiner als der Reihenabstand, verlagert den Abbruch auf die Nachbarstellung. Erst die Rundung $0{,}1$\,mm mit 12 Facetten rechnet alle geprüften Stellungen in sechs Inkrementen je Schritt ohne Warnung durch. Eine berührende konvexe Kante braucht demnach mindestens den Reihenabstand der Gegenfläche in Gleitrichtung als Radius. In der Serie brach mit dieser Kante noch eine von 396 Rechnungen ab.

Gegenüber der scharfen Kante verändert die Rundung das Ergebnis nur dort, wo die Kopfkante des Ritzels Last trägt. An 68 gemeinsamen Stellungen von Studie und Serie unter W1 steigt die Ritzeldrehung an B und E um $5{,}3$\,\% bis $6{,}8$\,\%. Der Kontaktdruck an der Kante fällt dort um $8{,}0$\,\% bis $10{,}8$\,\%, die Zugspannung $\sigma_{\mathrm{1}}$ um $4{,}7$\,\% bis $8{,}7$\,\%. An A, C und D bleiben alle Größen innerhalb $0{,}3$\,\%. Der Abschnitt des Wälzwegs um B und E, in dem die Rundung wirkt, ist etwa $0{,}3$\,mm lang. Die Größtwerte über dem ganzen Weg ändern sich bei $12$ und $16$\,Nm um höchstens $2{,}5$\,\%, bei $8$\,Nm um bis zu $4{,}9$\,\%.

Ob diese Empfindlichkeit an der Kante oder am Netz der Radflanke hängt, klärte ein Nachtlauf mit feinerer Flanke. Die Zahl der Reihen über der Zahnhöhe, der einzige Hebel für den Reihenabstand in Gleitrichtung, stieg von 18 auf 36 und 54. Der Abstand sank von etwa $0{,}1$\,mm auf etwa $0{,}05$ und $0{,}033$\,mm. Jedes Netz wurde mit Ritzelkanten von $0{,}1$, $0{,}05$ und $0{,}02$\,mm gepaart, W1 bei QS an den Normpunkten A bis E, die Kante $0{,}05$\,mm zusätzlich mit W2. \tabref{tab:netz_rundung} stellt die Änderungen gegen die Serie bei $12$\,Nm zusammen, dem Schritt, in dem alle Rechnungen durchliefen. Bei $16$\,Nm bricht W1 mit der Kante $0{,}05$\,mm auf 36 Reihen an B und E und auf 54 Reihen an E nach dem Muster der Studie ab. Mit $0{,}02$\,mm laufen alle fünf Stellungen auf beiden Netzen durch, ebenso jede Rechnung mit W2. Der Abbruch hängt damit am Verhältnis von Facettenlänge zu Reihenabstand und nicht am Radius allein.

\begin{table}[htbp]
    \caption{Netzeffekt gegen Rundungseffekt: Änderung der Werte an den Normpunkten gegen das Netz mit 18 Reihen über der Zahnhöhe und der Ritzelkante $0{,}1$\,mm, Ringkörper, W1, quasistatische Karte, $12$\,Nm. Drehung des Ritzels, größte Zugspannung $\sigma_{\mathrm{1}}$ der Fußrundung und größter Kontaktdruck $p$ unter der Kante, alle Angaben in \%.}
    \label{tab:netz_rundung}
    \centering
    \footnotesize
    \setlength{\tabcolsep}{3pt}
    \begin{tabular}{lllllllll}
        Reihen & Kante in mm & Drehung C & Drehung B & Drehung E & $\sigma_{\mathrm{1}}$ C & $\sigma_{\mathrm{1}}$ B & $\sigma_{\mathrm{1}}$ E & $p$ B     \\
        \hline
        36     & $0{,}1$     & $-1{,}6$  & $-2{,}0$  & $-2{,}2$  & $-1{,}7$                & $+3{,}7$                & $+0{,}9$                & $+50{,}6$ \\
        36     & $0{,}05$    & $-3{,}5$  & $-4{,}9$  & $-5{,}1$  & $-3{,}3$                & $+7{,}6$                & $+3{,}8$                & $+59{,}9$ \\
        36     & $0{,}02$    & $-3{,}5$  & $-5{,}0$  & $-5{,}3$  & $-3{,}3$                & $+9{,}3$                & $+5{,}9$                & $+59{,}9$ \\
        54     & $0{,}1$     & $-1{,}9$  & $-2{,}6$  & $-2{,}8$  & $-1{,}6$                & $+5{,}5$                & $+2{,}6$                & $+93{,}6$ \\
        54     & $0{,}05$    & $-3{,}8$  & $-4{,}8$  & $-5{,}0$  & $-3{,}9$                & $+8{,}1$                & $+4{,}1$                & $+97{,}0$ \\
        54     & $0{,}02$    & $-4{,}4$  & $-4{,}9$  & $-5{,}2$  & $-4{,}3$                & $+9{,}3$                & $+5{,}8$                & $+96{,}4$ \\
    \end{tabular}
\end{table}

% TK erledigt (Ritzeldrehung in 3.6 definiert: Verdrehwinkel bis zum Gleichgewicht aus Zahnbiegung und Kontakt, nicht das Flankenspiel; zum Kantendruck siehe den neuen Absatz am Ende): Drehung hatten wir glaube auch schonmal gesprochen, dass du eigentlich das Verdrehwinkelspiel meinst oder meinst du hier was anderes? Bitte nochmal besser spezifizieren, was du mit Drehung aussagen möchtest. Und wir haben ja jetzt mit 0,1 Rundung weitergearbeitet, aber ich finde es krass, dass bei 54 Reihen und 0,02er Rundung alles durchgelaufen ist und der Kontaktdruck +96 % hat. Das ist ja verrückt.

Der Netzeffekt allein zeigt sich mit der Kante $0{,}1$\,mm auf 36 und 54 statt 18 Reihen. Er senkt die Ritzeldrehung an A, C und D um bis zu $1{,}9$\,\% und an B und E um $1{,}8$\,\% bis $3{,}2$\,\%. An A, C und D bleibt die Zugspannung $\sigma_{\mathrm{1}}$ innerhalb $3$\,\%, an B steigt sie um $2{,}2$\,\% bis $5{,}5$\,\%. Der größte Kontaktdruck unter der Kante wächst dagegen mit jeder Verfeinerung weiter. An B steigt er mit 36 Reihen um $51$ bis $60$\,\% und mit 54 Reihen um $94$ bis $97$\,\% (\tabref{tab:netz_rundung}), an E in ähnlicher Höhe. Ein konvergierter Wert existiert nicht. Mit der aufgelösten scharfen Kante, $0{,}02$\,mm auf 54 Reihen, dreht das Ritzel gegenüber der Serie um $2{,}0$\,\% bis $4{,}4$\,\% weniger an A, C und D. An B und E sind es $4{,}5$\,\% bis $5{,}9$\,\% weniger. Die Zugspannung liegt an C um $2{,}4$\,\% bis $4{,}3$\,\% tiefer, an B um $7{,}5$\,\% bis $9{,}4$\,\% und an E um $4{,}7$\,\% bis $6{,}0$\,\% höher.

% TK erledigt (Text auf die Werte an B der Tabelle gebracht, E nur qualitativ): in der Tabelle stehen keine 103 % sondern 96,4 %. Bitte Zahlen im Text passend zu den Tabellen oder irgendwas ist durcheinander gekommen.

Die reale Kante des Prüfstandsritzels 95912 wurde an drei Zähnen im Konturscan aufgenommen. An jede Kante wurden ein tangentialer Bogen vom Radius $\rho$ und eine symmetrische Fase mit dem Schenkel $c$ gefittet. \figref{fig:kopfkante_95912} zeigt Zahn 18 mit Bogenradien von $0{,}004$ bis $0{,}008$\,mm bei einer Restabweichung von $0{,}9$ bis $1{,}3$\,µm. Zahn 1 liegt im Mittel bei $0{,}004$\,mm. Mikroskopaufnahmen der drei Zahnköpfe zeigen die gedrehte Kopffläche und die gewälzfräste Flanke an einer Kante ohne Fasenfacette und ohne sichtbares Rundungsband. Daraus folgt nur ein Radius unter etwa $20$\,µm. Auch das ältere Ritzel 86481 trägt mit einem Bogen von im Mittel $0{,}008$\,mm eine scharfe Kante.

\begin{figure}[htbp]
    \centering
    \includegraphics[width=0.8\textwidth]{kap4/kopfkante_95912}
    \caption{Kopfkante des Prüfstandsritzels 95912 im Konturscan, Zahn 18, rechte und linke Kante des ersten Zahnkopfs: Messpunkte, Lage der scharfen Ecke, tangentialer Bogen und symmetrische Fase als Fit, die Werte beider Formen stehen im Text.}
    \label{fig:kopfkante_95912}
\end{figure}

% TK (folgt mit Paket B4, Bild wird neu gezeichnet): Bitte Kopf 2 rausnehmen. Da ist eh was falsch mit einer Geraden als Fasenfit und Kopf 1 rechts und links reicht.

Die Rundung $0{,}1$\,mm mit 12 Facetten ist damit eine Modellwahl und keine Eigenschaft des Teils. Dass trotz der Verdopplung des Kantendrucks mit dieser Rundung weitergerechnet wurde, hat drei Gründe. Der Druck unter einer nahezu scharfen Kante ist ein Linienkontakt ohne konvergierenden Wert, jede Verfeinerung hebt die Spitze weiter. Er ist deshalb keine Bemessungsgröße dieser Arbeit und wird in den Ergebniskapiteln als netz- und kantenabhängige Größe geführt. Die Zielgrößen am Fuß ändern sich mit Kante und Reihenzahl an B und E um höchstens $9{,}4$\,\% in $\sigma_{\mathrm{1}}$ und rund $5$\,\% in der Ritzeldrehung, an den übrigen Normpunkten und an den Größtwerten um wenige Prozent. Und die reale Kante mit einem Radius unter $20$\,µm ist schärfer als jedes Modell. Sie arbeitet sich im Versuch in die Radflanke ein (\secref{sec:schadensbilder}), so dass der Kantendruck des unverschlissenen Modells das Teil ohnehin nicht beschreibt. Ein Modell mit scharfer Kante bräuchte $0{,}02$\,mm auf mindestens 36 Reihen.

% TK erledigt (drei Gründe im Text; Mapping gab es nur für das 18-Reihen-Netz, die 36/54-Piloten liefen mit W1/W2): Wenn das im Kontaktdrcuk Abweichungen von 100 % waren grob, dann bitte hier noch eine bessere erklärung, warum doch einfach mit der 0,1er Rundung weiter gearbeitet wurde, da wir ja das mapping auf die feineren Reihen durchgeführt hatten. 

% Herkunft:
% Absatz 1 (Studie 60 je Teilung): ADR-116 Punkt 36 (architecture_decisions.md, Amendment 2026-10-09);
%   Teilgitter-Zahlen aus 80_output/fe/kst_e/study_20layers_bore33_60perpitch_gerechnet/path_report.md
%   (LOAD_16NM: 30 je Teilung s1_max −0,14 %, 20 je Teilung s3_min −1,59 % und Versatz ≤ 0,035 mm
%   ausser p_max +0,117 mm bei −0,03 %, 15 je Teilung s3_min −2,12 %, 10 je Teilung s3_min −2,47 %;
%   LOAD_8NM: 30 je Teilung s3_min −1,74 %, 20 je Teilung 0,00 %).
% Absatz 2 (Serie, Rand): series_..._30perpitch_margin1_qs_gerechnet/analyse/analyse.md Kopf (396 Dateien,
%   30 je Teilung, Rand 1,00, 99 Stellungen je Stufe) und Abschnitt 6 (15 je Teilung +2,9 % W3 16 Nm, 10 je Teilung +5,5 %);
%   ADR-116 Punkt 41 (Rand 0,5 Teilung zu kurz, 10 und 30 % der Zahnkraft) und Punkt 47 (a)
%   (0,47 bis 0,60 vor A, 0,35 bis 0,48 hinter E).
% Absatz 3 (Abbrüche, Rundung 0,1 mm mit 12 Facetten): ADR-116 Punkte 37 und 43; known_limits.md FE-19
%   (zwei Arten, Versatz 0,025 mm, Zweige der Normalen, Regel Reihenabstand, 1 von 396 pos_073);
%   series analyse.md Abschnitt 1 (pos_073_W1_QS bei 16 Nm, 3,838 mm ab A, r 25,800 mm).
% Absatz 4 (scharf gegen gerundet): series analyse.md Abschnitt 5 (68 Stellungen, Tabellen 8/12/16 Nm:
%   Drehung B +6,7/+5,6/+5,3, E +6,8/+5,9/+5,4; Druck B −9,9/−10,8/−9,1, E −9,0/−9,5/−8,0;
%   σ1 B −7,5/−8,3/−8,7, E −4,7/−5,3/−5,7; A, C, D höchstens +0,3 % (Druck an C);
%   Maxima 8 Nm bis +4,9 %, 12 Nm bis +1,6 %, 16 Nm bis +2,5 %); ADR-116 Punkt 47 (c) (Fenster 0,3 mm);
%   known_limits.md FE-17. Abweichung: ADR nennt "Maxima innerhalb 2,5 %", die Tabelle bei 8 Nm +4,6 % (σ1)
%   und +4,9 % (|σ3|); der Text folgt der Tabelle.
% Absatz 5 und Tabelle (Nachtlauf, Netz gegen Rundung): ADR-116 Punkt 50 mit Amendment und Outcome
%   (18/36/54 Reihen, Abstände, Abbrüche); fe_simulation_notes.yaml outcome der sechs Ordner
%   pilot_20layers_bore33_hoehe{36,54}_round0p{02,05,1}f12_*_qs_AtoE_gerechnet; Tabelle: Abschnitt 5
%   (Vergleich mit der Serie) der sechs analyse.md, Block 12 Nm, Spalten C, B, E.
% Absatz 6 (Netzeffekt, aufgelöste Kante): dieselben sechs analyse.md, Abschnitt 5, Blöcke 8/12/16 Nm,
%   Spannen über die drei Schritte selbst gebildet (Netzeffekt 0,1 mm: Drehung A/C/D +0,1 bis −1,9 %,
%   B/E −1,8 bis −3,2 %, σ1 A/C/D −3,0 bis +0,2 %, B +2,2 bis +5,5 %; Druck B/E +48,3 bis +57,2 % (36)
%   und +86,8 bis +102,8 % (54); aufgelöste Kante hoehe54_round0p02: Drehung A/C/D −2,0 bis −4,4 %,
%   B/E −4,5 bis −5,9 %, σ1 C −2,4 bis −4,3 %, B +7,5 bis +9,4 %, E +4,7 bis +6,0 %); W2 gegen W1
%   Kontaktdruck an B/E (−32,7/−45,0 % auf 36 Reihen, −37,5/−51,3/−58,3 % auf 54; Abschnitt 3 der
%   round0p05-Ordner) aus dem Text gestrichen, Stoff für 4.3. Abweichung: ADR-116 Punkt 50 nennt σ1 an A, C, D "innerhalb ±1 %" und
%   Drehung A/C/D "−2,6 bis −3,5 %"; die 12-Nm- und 16-Nm-Blöcke reichen weiter (σ1 A/D 16 Nm −3,0/−2,8 %,
%   Drehung C 12 Nm −4,4 %); der Text folgt den Tabellen.
% Absatz 7 und Bild (reale Kante): 80_output/messungen/konturscan/pinion/bericht.md Kopf (Fit Bogen/Fase),
%   Tabelle je Scan (95912_Z01 ρ 0,004, Z18 ρ 0,006; Z36 ρ 0,091 wegen eines als Kopffläche erkannten
%   Laufs 0,28 mm unter den anderen nicht genannt, ADR-117 Punkt 4) und
%   je Teil (86481 ρ Mittel 0,008, max 0,018); SVG kopfkante_95912_95912_Z18 (ρ 0,008/0,004/0,005 mm,
%   RMS 1,3/0,9/1,3 µm); ADR-117 Amendment 2026-10-12 Punkt 4 (Maschinenhöhe 87,57 bis 87,62 mm,
%   Bogen 0,002 bis 0,008 mm, Mikroskop F4092 Z1/Z18/Z36, Maßbalken 200 µm, "unter etwa 20 µm",
%   Rücknahme 25 bis 27 µm gegen fKo −22,7/−23,4 µm). Hinweis: Die Vorgabe "Kopfrücknahme nur beim
%   älteren Ritzel 86481" steht gegen bericht.md (Gerade am Nennkopfkreis 86481 20,8 µm, 95912 25,9 µm)
%   und ADR-117; deshalb der Platzhalter. Die Fasen-Schenkel der bericht.md-Tabelle (0,062/0,093 mm)
%   weichen von ADR-117 (0,001 bis 0,004 mm) und vom SVG (0,004/0,002/0,274 mm) ab; der Text nennt
%   deshalb nur die Bogenradien.
% Absatz 8 (Folgerung): ADR-116 Punkt 50 Outcome (Decision) und Punkt 43; known_limits.md FE-17 (Modellparameter).
% Bildwunsch: 80_output/fe/kst_e/series_20layers_bore33_round0p1f12_30perpitch_margin1_qs_gerechnet/analyse/vergleich_referenz_LOAD_16NM.svg,
%   Rundung 0,1 mm gegen scharfe Kante über dem Eingriffsweg (Fenster um B und E; Legende ohne "Serie"/"Referenz").
% Bildwunsch: Kurven der 60-je-Teilung-Studie mit den Teilgittern 30/20/15/10 (im Studienordner liegen nur mesh_*.png,
%   bilder_V2/study_* nicht gefunden; neu zu zeichnen aus den Stellungsdateien des Studienordners).
% Bildwunsch: 80_output/fe/kst_e/pilot_20layers_bore33_hoehe54_round0p02f12_w1_qs_AtoE_gerechnet/analyse/vergleich_referenz_LOAD_16NM.svg
%   (oder ein Sammelbild der sechs Nachtlauf-Ordner als Marker an A bis E), Netzeffekt gegen Rundungseffekt.
% Bildwunsch: Mikroskopaufnahme F4092 95912 Z18 (Inventar, Art microscope) als Ergänzung zum Konturscan der Kante.
